\documentclass[11pt]{article}
\usepackage[margin=1in]{geometry}
\usepackage{amsmath,amssymb,amsthm}
\usepackage{graphicx}
\usepackage[hidelinks]{hyperref}

\newtheorem{theorem}{Theorem}
\newtheorem{lemma}{Lemma}
\newtheorem{remark}{Remark}

\DeclareMathOperator{\Rg}{Rg}
\DeclareMathOperator{\Ree}{Re}

\begin{document}

\begin{center}
{\Large Normal Hilbert semigroups and the fractional-range stability condition}\\[0.4em]
{\small Run \texttt{fa\_banach\_001}, arXiv:math/0512460, partial answer}
\end{center}

\section*{Source problem}

Borichev, Chill, and Tomilov prove in \cite{BCT} that, for a bounded
\(C_0\)-semigroup generated by \(A\) on a Banach space with Fourier type
\(p\in(1,2]\), density of the set
\[
M_p(A):=\left\{x:\ \lim_{\alpha\to0+}
\alpha^{(p-1)/p}R(\alpha+i\beta,A)x=0
\text{ for every }\beta\in\mathbb R\right\}
\]
implies strong stability.  They also prove the algebraic sufficient condition
\[
\bigcap_{\beta\in\mathbb R}\Rg(i\beta-A)^{1/p}\ \text{ is dense.}
\]
Remark 5.11 asks whether these conditions characterize stability in the
Hilbert case \(p=2\).

\begin{center}
\includegraphics[width=0.9\linewidth]{figures/open_problem_crop.png}
\end{center}

\section*{Partial answer}

\begin{theorem}\label{thm:normal}
Let \(A\) be a normal generator of a bounded \(C_0\)-semigroup
\((T(t))_{t\ge0}\) on a complex Hilbert space \(H\).  Then the following are
equivalent:
\begin{enumerate}
\item \((T(t))_{t\ge0}\) is strongly stable;
\item \(M_2(A)\) is dense in \(H\);
\item \(\displaystyle \bigcap_{\beta\in\mathbb R}\Rg(i\beta-A)^{1/2}\) is
dense in \(H\).
\end{enumerate}
\end{theorem}

\begin{proof}
The implications (2) \(\Rightarrow\) (1) and (3) \(\Rightarrow\) (1) are
Theorem 5.7 and Corollary 5.10 of \cite{BCT}, specialized to \(p=2\).
It remains to prove the converse under normality.

Let \(E\) denote the spectral measure of \(A\).  Since \(A\) generates a
bounded normal semigroup, the spectral support is contained in the closed left
half-plane and
\[
T(t)=e^{tA}=\int e^{tz}\,dE(z).
\]
Strong stability is equivalent to \(E(i\mathbb R)=0\).  Indeed, if
\(E(i\mathbb R)x=x\), then \(\|T(t)x\|=\|x\|\) for all \(t\ge0\), so strong
stability forces \(x=0\).  Conversely, if \(E(i\mathbb R)=0\), then
\(|e^{tz}|\le1\) and \(e^{tz}\to0\) for \(E\)-almost every spectral point, so
dominated convergence gives \(\|T(t)x\|\to0\) for every \(x\in H\).

Assume now that the semigroup is strongly stable.  For \(\varepsilon>0\), set
\[
H_\varepsilon:=E\{z:\Ree z\le -\varepsilon\}H .
\]
Because \(E(i\mathbb R)=0\), the union \(\bigcup_{\varepsilon>0}H_\varepsilon\)
is dense in \(H\).  Fix \(x\in H_\varepsilon\) and \(\beta\in\mathbb R\).  On
the spectral support of \(x\),
\[
|i\beta-z|\ge \varepsilon .
\]
Hence the multiplier \((i\beta-z)^{-1/2}\) is bounded there.  By the normal
functional calculus,
\[
y_\beta:=\int (i\beta-z)^{-1/2}\,dE(z)x
\]
belongs to \(H\), and
\[
(i\beta-A)^{1/2}y_\beta=x.
\]
Thus \(x\in\Rg(i\beta-A)^{1/2}\) for every \(\beta\), and so
\[
\bigcup_{\varepsilon>0}H_\varepsilon
\subset
\bigcap_{\beta\in\mathbb R}\Rg(i\beta-A)^{1/2}.
\]
The latter intersection is therefore dense.

The same spectral gap proves density of \(M_2(A)\).  If \(x\in H_\varepsilon\),
then for every \(\beta\in\mathbb R\),
\[
\left\|\alpha^{1/2}R(\alpha+i\beta,A)x\right\|
\le
\frac{\alpha^{1/2}}{\varepsilon}\|x\|\longrightarrow0
\qquad(\alpha\to0+).
\]
Thus \(\bigcup_{\varepsilon>0}H_\varepsilon\subset M_2(A)\), and \(M_2(A)\)
is dense.
\end{proof}

\begin{remark}
The normality hypothesis is essential for this packet's proof.  The argument
uses the spectral theorem to manufacture a common dense subspace lying a
positive distance from the entire imaginary axis.  It does not address the
source problem for arbitrary Hilbert-space generators.
\end{remark}

\section*{Ledger note}

Cheap run indexes were searched for arXiv:0512460 and the core stability
keywords; no duplicate packet was found.  A web search on 2026-07-03 for the
exact source title and the Hilbert/Fourier-type stability question found the
source paper but no later public settlement of the full problem.

\begin{thebibliography}{9}
\bibitem{BCT}
A. Borichev, R. Chill, and Y. Tomilov,
\emph{Uniqueness theorems for (sub-)harmonic functions with applications to
operator theory}, arXiv:math/0512460, 2005.
\end{thebibliography}

\end{document}

